% from ega4-2-fr.tex
\oldpage[IV-2]{132}

Il est clair en effet que
$\varphi_{\beta\alpha}(\mathfrak N_\alpha)\subset\mathfrak N_\beta$ pour
$\alpha\leq\beta$, et que $\varinjlim\mathfrak N_\alpha$ est un idéal
nilpotent de $A$, donc contenu dans le nilradical $\mathfrak N$ de $A$.
Inversement, si $x\in\mathfrak N$, on a $x^h=0$ pour un entier $h$; on peut
écrire $x=\varphi_\alpha(x_\alpha)$ pour un indice $\alpha$ et un
$x_\alpha\in A_\alpha$; la relation $\varphi_\alpha(x_\alpha^h)=0$ entraîne
l'existence d'un $\beta\geq\alpha$ tel que
$\varphi_{\beta\alpha}(x_\alpha^h)=0$, ou encore $x_\beta^h=0$ en posant
$x_\beta=\varphi_{\beta\alpha}(x_\alpha)$; par suite
$x_\beta\in\mathfrak N_\beta$ et comme $x=\varphi_\beta(x_\beta)$ cela achève
de prouver le corollaire, compte tenu de (5.13.1).

\begin{proposition}[5.13.3]
\label{IV.2.s5.13.3-fr}
\begin{enumerate}
\item[(i)] Pour tout indice $\alpha$, soit $\mathfrak p_\alpha$ un idéal
premier de $A_\alpha$ tel que, pour $\alpha\leq\beta$, on ait
$\varphi_{\beta\alpha}(\mathfrak p_\alpha)\subset\mathfrak p_\beta$. Alors
$\mathfrak p=\varinjlim\mathfrak p_\alpha$ est un idéal premier de $A$. En
particulier, si les $A_\alpha$ sont intègres il en est de même de $A$.
\item[(ii)] Si de plus on a
$\mathfrak p_\alpha=\varphi_{\beta\alpha}^{-1}(\mathfrak p_\beta)$ pour
$\alpha\leq\beta$, alors $A_{\mathfrak p}$ s'identifie canoniquement à
$\varinjlim(A_\alpha)_{\mathfrak p_\alpha}$.
\item[(iii)] Si $A$ est un anneau local, $\mathfrak p$ son idéal maximal, et
si les homomorphismes $\varphi_\alpha$ sont injectifs, alors, en posant
$\mathfrak p_\alpha=\varphi_\alpha^{-1}(\mathfrak p)$, on a
$A=\varinjlim(A_\alpha)_{\mathfrak p_\alpha}$ et les homomorphismes
$(A_\alpha)_{\mathfrak p_\alpha}\to A$ sont injectifs.
\end{enumerate}
\end{proposition}

Les deux assertions de (i) sont en fait équivalentes, puisque
$A/\mathfrak p=\varinjlim(A_\alpha/\mathfrak p_\alpha)$. Supposons donc tous
les $A_\alpha$ intègres, et soient $x$, $y$ deux éléments de $A$ tels que
$xy=0$. Il existe donc un indice $\alpha$ tel que
$x=\varphi_\alpha(x_\alpha)$, $y=\varphi_\alpha(y_\alpha)$ avec
$x_\alpha,y_\alpha$ dans $A_\alpha$, et
$\varphi_\alpha(x_\alpha y_\alpha)=0$; mais cela entraîne l'existence d'un
indice $\beta\geq\alpha$ tel que, si l'on pose
$x_\beta=\varphi_{\beta\alpha}(x_\alpha)$,
$y_\beta=\varphi_{\beta\alpha}(y_\alpha)$, on ait $x_\beta y_\beta=0$; comme
$A_\beta$ est intègre, on a $x_\beta=0$ ou $y_\beta=0$ et comme
$x=\varphi_\beta(x_\beta)$ et $y=\varphi_\beta(y_\beta)$, cela achève de
montrer que $A$ est intègre.

Prouvons maintenant (ii). L'hypothèse entraîne que les
$A_\alpha-\mathfrak p_\alpha$ forment un système inductif de parties
multiplicatives des $A_\alpha$, dont $A-\mathfrak p$ est la limite inductive;
par suite les anneaux locaux $(A_\alpha)_{\mathfrak p_\alpha}$ forment un
système inductif (où les homomorphismes de transition sont locaux), et les
homomorphismes canoniques
$(A_\alpha)_{\mathfrak p_\alpha}\to A_{\mathfrak p}$ un système inductif
d'homomorphismes, dont la limite inductive
\[
\psi:\varinjlim (A_\alpha)_{\mathfrak p_\alpha}
\longrightarrow A_{\mathfrak p}
\]
est un homomorphisme surjectif. Reste à voir que $\psi$ est injectif. Or, si
$x_\alpha\in A_\alpha$ et
$s_\alpha\in A_\alpha-\mathfrak p_\alpha$ sont tels que
$\varphi_\alpha(x_\alpha)/\varphi_\alpha(s_\alpha)=0$ dans $A_{\mathfrak p}$,
il existe $s'\in A-\mathfrak p$ tel que
$s'\varphi_\alpha(x_\alpha)=0$, et, en remplaçant au besoin $\alpha$ par un
indice plus grand, on peut supposer que $s'=\varphi_\alpha(s'_\alpha)$ avec
$s'_\alpha\in A_\alpha-\mathfrak p_\alpha$, d'où
$\varphi_\alpha(s'_\alpha x_\alpha)=0$; il y a donc un $\beta\geq\alpha$ tel
que $\varphi_{\beta\alpha}(s'_\alpha x_\alpha)=0$ dans $A_\beta$, ce qui
entraîne
$\varphi_{\beta\alpha}(x_\alpha)/\varphi_{\beta\alpha}(s_\alpha)=0$ dans
$(A_\beta)_{\mathfrak p_\beta}$ et prouve donc que l'image canonique de
$x_\alpha/s_\alpha$ dans
$\varinjlim(A_\alpha)_{\mathfrak p_\alpha}$ est nulle.

Enfin, la première assertion de (iii) résulte de (ii) et de ce que
$A_{\mathfrak p}=A$. En outre, les éléments de
$A_\alpha-\mathfrak p_\alpha$ sont inversibles dans $A$, donc on a aussi
$A_{\mathfrak p_\alpha}=A$, et le fait que
$(A_\alpha)_{\mathfrak p_\alpha}\to A$ soit injectif se déduit par platitude
de l'hypothèse que $\varphi_\alpha$ est injectif.

\begin{corollary}[5.13.4]
\label{IV.2.s5.13.4-fr}
Supposons que les $A_\alpha$ soient tous intègres, et que les
$\varphi_{\beta\alpha}$ soient injectives. Alors, si $A'_\alpha$ est la clôture
intégrale de $A_\alpha$, $A'=\varinjlim A'_\alpha$ est la clôture intégrale de
$A$. En particulier, si $A_\alpha$ est intégralement clos pour tout $\alpha$,
il en est de même de $A$. Si les $A_\alpha$ sont des anneaux locaux unibranches
(resp. géométriquement unibranches) (0, 23.2.1) et les
$\varphi_{\beta\alpha}$ des homomorphismes locaux injectifs, $A$ est unibranche
(resp. géométriquement unibranche).
\end{corollary}

Si $K_\alpha$ est le corps des fractions de $A_\alpha$, il résulte de (5.13.3,
(ii)) appliqué aux $\mathfrak p_\alpha=0$, que
$K=\varinjlim K_\alpha$ est le corps des fractions de $A$, et
$A'=\varinjlim A'_\alpha$ s'identifie à

